% !TeX spellcheck = en_US
\documentclass[a4paper,12pt]{article}

\usepackage{fullpage}
\usepackage[utf8]{inputenc}
\usepackage{fourier}
\usepackage{amsmath}
\usepackage{color}
\usepackage{graphicx}
\usepackage{titlesec}

\titleformat{\subsection}[hang]{\large\bfseries}{\alph{subsection})\quad}{0pt}{}

\newcommand{\twodo}[1]{\textcolor{red}{\textbf{todo:} #1}}

\title{\textbf{Exercises for Image Processing 1}\\Problem Sheet 4}
\author{Tim Dobert\\6427948 \and Konstantin M\"ollers\\6313136}

\begin{document}
	\maketitle	
	
	\section{Theoretical Problems}
	\subsection{Histograms and Noise}
	
	\begin{itemize}
		\item Ein Histogramm repräsentiert, wie viele Pixel einen bestimmten Grauwert haben
		\item Das Histogramm hätte zwei Spitzen, eine für das Fließband, die andere für die Pakete
		\item Wenn das Rauschen Gaussverteilt ist, sollte sich das Histogramm eines normalen Bildes nicht großartig ändern.
		\item Für jeden Pixel der dunkler wird, wird ein anderer Heller.
		\item Außnahme ist ein komplett einfarbiges Bild
	\end{itemize}
	
	\subsection{Projections}
	\subsection{Filters}
	\subsection{Convolution}
	
	If you look at B1, you can easily imagine that if we represent it by a function it would look like a rectangular peak. If you convolute two rectangular peaks with each other, they look like a triangular peak. So now take this one dimension higher, you can imagine that effect taking place in each direction of the image and so you get a rectangular shaped cone. You can see that resulting image B2 below.
	
	\noindent
	\begin{tabular}{@{}p{.45\linewidth}@{\hspace*{.1\linewidth}}p{.45\linewidth}@{}}
%		\begin{figure}
%			\centering
			\includegraphics[width=\linewidth]{B1}
%			\caption{Input images B1}
%		\end{figure}
		&
%		\begin{figure}
%			\centering
			\includegraphics[width=\linewidth]{B2}
%			\caption{Input images B1}
%		\end{figure}
\\
	Input images \textit{B1} &
	Resulting convoluted image \textit{B2}
	\end{tabular}
	
	
	\section{Practical Problems}
	\subsection{Gray Value Normalization}
	\subsection{Fourier-Transform}
	
	1D DFT:
	
	\begin{align*}
		\hat{a}_l &= \sum\limits_{n = 0}^{N - 1} a_n \cdot e^{-2 \pi i \frac{nl}{N}} \\
		\hat{a}_k &= \sum\limits_{m = 0}^{M - 1} a_m \cdot e^{-2 \pi i \frac{mk}{M}} \\
	\end{align*}
	
	2D DFT:
	
	\begin{align*}
	\hat{a}_{kl} &= \sum\limits_{m = 0}^{M - 1} \sum\limits_{n = 0}^{N - 1} a_{m, n} \cdot e^{-2 \pi i \frac{mk}{M}} \cdot e^{-2 \pi i \frac{nl}{N}} \\
	\end{align*}
\end{document}
